\documentclass[11pt,a4paper]{article}
\usepackage[T1]{fontenc}
\usepackage[utf8]{inputenc}
\usepackage{lmodern,microtype}
\usepackage[margin=27mm]{geometry}
\usepackage{amsmath,amssymb,amsthm,mathtools}
\usepackage{booktabs,array,enumitem,xcolor}
\usepackage[colorlinks=true,linkcolor=blue!45!black,citecolor=blue!45!black,urlcolor=blue!45!black]{hyperref}
\usepackage{fancyhdr}
\pagestyle{fancy}\fancyhf{}
\fancyhead[L]{\small Antisymmetric spectral invariants}
\fancyhead[R]{\small T71--T75}
\fancyfoot[C]{\thepage}
\setlength{\headheight}{14pt}
\setlength{\parskip}{4pt}
\setlength{\parindent}{1.1em}
\setlist{itemsep=3pt,topsep=5pt}
\newtheorem{theorem}{Theorem}
\renewcommand{\thetheorem}{T\arabic{theorem}}
\setcounter{theorem}{70}
\newtheorem{lemma}{Lemma}[section]
\newtheorem{proposition}[lemma]{Proposition}
\newtheorem{corollary}[lemma]{Corollary}
\theoremstyle{definition}
\newtheorem{problem}{Auxiliary problem}
\theoremstyle{remark}
\newtheorem*{remark}{Remark}
\newcommand{\C}{\mathbb C}
\newcommand{\Z}{\mathbb Z}
\newcommand{\A}{\mathcal A}
\newcommand{\X}{\mathcal X}
\newcommand{\gr}{\operatorname{gr}}
\newcommand{\rank}{\operatorname{rank}}
\newcommand{\tr}{\operatorname{tr}}
\newcommand{\Spec}{\operatorname{Spec}}
\newcommand{\length}{\operatorname{length}}
\newcommand{\rev}{\operatorname{rev}}
\newcommand{\ini}{\operatorname{in}}
\DeclareMathOperator{\diag}{diag}
\title{Antisymmetric Floquet-Isospectral Potentials:\\
Existence, Local Multiplicity, and Singular Period Inflations}
\author{KEIO}
\date{24 October 2026\\[3pt]\small Research manuscript: T71--T75}
\begin{document}
\maketitle
\begin{abstract}
For the discrete periodic Schr\"odinger operator on the integer lattice, we study
the finite scheme of complex potentials that are antisymmetric under reversal
and Floquet isospectral to zero. If the period is $q$ and $m=\lfloor q/2\rfloor$,
we identify the associated graded local algebra at zero with
$\C[x_1,\ldots,x_m]/(x_1^2,\ldots,x_m^2)$.
Its length is therefore $2^m$. Comparing this local length with the known global
length $2^m m!$ proves the existence of a nonzero antisymmetric potential for
every $q>3$, answering Conjecture~1 of Faust, Friedman, O'Malley, Ramos, and
Sharma. We also prove an exact differential-rank identity for period repetition.
It produces a nonzero singular point at period~15 and inequivalent nonzero
singular points at period~20, disproving the odd regularity and even uniqueness
assertions of their Conjecture~3. The same argument gives arbitrarily many
singular symmetry classes. The five theorem labels organize these results;
they represent resolutions of two previously numbered conjectures, not five
independent open problems. Historical priority is discussed separately from
the proofs.
\end{abstract}

\section{Targets, contribution, and the auxiliary questions}

The prior targets are Conjectures~1 and~3 of \cite{FFORS}, a preprint submitted
on 9 February 2026. In the notation introduced below, their assertions are:
\begin{enumerate}[label=(\roman*)]
\item For every integer $q>3$, the antisymmetric zero-isospectral scheme
      $\X_q$ has a nonzero point (Conjecture~1).
\item At odd $q$, all nonzero points are regular; at even $q>6$, there is
      one nonzero singular symmetry class, with multiplicity $2^{q/2-2}$;
      the origin has multiplicity $2^{\lfloor q/2\rfloor}$ (Conjecture~3).
\end{enumerate}
Here multiplicity and regularity refer to the defining spectral ideal, with
its possible nilpotents. Also, a $q\Z$-periodic potential need not have least
period $q$. These conventions matter for the counterexamples.

We introduce two auxiliary questions as the entry points to the argument.
They are research questions formulated here, without an assertion that the
questions themselves have never previously been considered.
\begin{problem}[The quadratic local-algebra question]\label{prob:local}
Can the lowest nonzero terms of the antisymmetric spectral equations be
simultaneously identified with independent squares? If so, does this determine
the entire associated graded local algebra at the zero potential?
\end{problem}
\begin{problem}[The repeated-point differential question]\label{prob:repeat}
If a potential of period $r$ is repeated to period $q=kr$, what happens to the
rank of the differential of the antisymmetric spectral map?
\end{problem}

Theorem~T71 answers Problem~\ref{prob:local}; Theorem~T72 applies it to the
unchanged Conjecture~1. Theorem~T73 answers Problem~\ref{prob:repeat};
Theorems~T74--T75 apply it to the unchanged Conjecture~3. The method connects a
local calculation to a global existence statement: cyclic correlations become
Fourier squares, the squares determine local length, and the difference between
global and local length forces nonzero solutions.

The transfer identities, Fourier transform, intersection theory, and regular
sequence arguments used here are classical. In particular, repetition preserves
Floquet isospectrality in prior work \cite[Section~4.2, equation~(12)]{CFK}.
Our proposed contribution is the exact local-algebra identification and the
resulting conjecture resolutions, together with the differential-rank
calculation. We do not assert that these basic tools are a historically new
general principle. Section~\ref{sec:scope} records the scope of the priority check.

\section{The precise spectral scheme}

Fix $q\geq3$ and $m=\lfloor q/2\rfloor$. Throughout, indices of a $q$-tuple
run from $0$ to $q-1$, with cyclic indices when indicated. Define
\[
 \A_q=\{v\in\C^q:v_j=-v_{q-1-j}\}.
\]
It has coordinates $u=(u_1,\ldots,u_m)$ through
\[
v=\begin{cases}
(u_1,\ldots,u_m,-u_m,\ldots,-u_1),&q=2m,\\
(u_1,\ldots,u_m,0,-u_m,\ldots,-u_1),&q=2m+1.
\end{cases}
\]
For $z\ne0$, let $L_v(z)$ be the cyclic tridiagonal matrix with diagonal $v$,
ones immediately above and below the diagonal, and corner entries
$(L_v(z))_{0,q-1}=z^{-1}$ and $(L_v(z))_{q-1,0}=z$.
The potentials considered here are complex; no nonzero real potential is
being asserted.

Set
\begin{equation}\label{eq:transfer}
 A_a(\lambda)=\begin{pmatrix}\lambda-a&-1\\1&0\end{pmatrix},\qquad
 M_q(v;\lambda)=A_{v_{q-1}}(\lambda)\cdots A_{v_0}(\lambda),\qquad
 D_q(v;\lambda)=\tr M_q(v;\lambda).
\end{equation}
The recurrence $f_{j+1}=(\lambda-v_j)f_j-f_{j-1}$ shows that the Floquet
condition is the condition that $z$ be an eigenvalue of $M_q$. Since
$\det M_q=1$, the usual monic characteristic-polynomial identity is
\begin{equation}\label{eq:floquet}
 \det(\lambda I-L_v(z))=D_q(v;\lambda)-z-z^{-1}.
\end{equation}
It can also be checked directly by determinant expansion. Consequently
$v$ is Floquet isospectral to zero precisely when
\begin{equation}\label{eq:zeroiso}
 D_q(v;\lambda)=D_q(0;\lambda)=2T_q(\lambda/2).
\end{equation}
Here $T_q$ is the Chebyshev polynomial of the first kind. The last equality
follows from Cayley--Hamilton for $A_0$, whose trace is $\lambda$ and determinant
is one.

The discriminant is invariant under cyclic permutation and reversal of the
potential. Cyclic invariance follows from the trace; reversal follows from
\eqref{eq:floquet} and reversal of the Floquet matrix. In addition, conjugating
each transfer matrix by $\diag(1,-1)$ gives
\[
 D_q(-v;-\lambda)=(-1)^qD_q(v;\lambda).
\]
On $\A_q$, reversal equals negation. Thus $D_q(v;\lambda)$ has parity $q$
in $\lambda$, and each coefficient is even as a polynomial in $u$.
The leading monic term is independent of $u$. Therefore
\begin{equation}\label{eq:coeffs}
 F_q(u;\lambda):=D_q(v;\lambda)-D_q(0;\lambda)
     =\sum_{\ell=1}^m f_{q,\ell}(u)\lambda^{q-2\ell}.
\end{equation}
Every $f_{q,\ell}$ has zero constant and linear terms. We use the schemes
\begin{equation}\label{eq:scheme}
 I_q=(f_{q,1},\ldots,f_{q,m})\subset\C[u_1,\ldots,u_m],\qquad
 \X_q=\Spec\bigl(\C[u_1,\ldots,u_m]/I_q\bigr).
\end{equation}
These generators differ only by signs from the specialized spectral
generators in \cite{FFORS}, so this is the same ideal and the same scheme.
Write $\mu_q(v)=\dim_\C\mathcal O_{\X_q,v}$ for its local length.

\begin{proposition}[Global length; also \cite{FFORS}, Theorem~4.6]\label{prop:global}
The scheme $\X_q$ is a finite complete intersection of total length
\[
 \length(\X_q)=2^m m!.
\]
\end{proposition}
\begin{proof}
The degree of $f_{q,\ell}$ in $u$ is $2\ell$, and its highest homogeneous
part is $(-1)^\ell e_\ell(u_1^2,\ldots,u_m^2)$. Indeed, the terms of highest
potential degree in the relevant characteristic coefficient come from the
diagonal product
\[
 \prod_{j=0}^{q-1}(\lambda-v_j)
  =\lambda^{q-2m}\prod_{j=1}^m(\lambda^2-u_j^2).
\]
If all $e_\ell(u_1^2,\ldots,u_m^2)$ vanish, then
$\prod_j(t-u_j^2)=t^m$, so all $u_j=0$.
The homogenizations of the $m$ equations therefore have no common projective
zero on the hyperplane at infinity. Their intersection is zero dimensional:
a positive dimensional projective component would meet that hyperplane.
Since there are $m$ hypersurfaces in $\mathbb P^m$ and their intersection
is zero dimensional, they form a complete intersection. Projective B\'ezout
gives length $\prod_{\ell=1}^m(2\ell)=2^m m!$. All its points lie in the
affine chart, proving the assertion.
\end{proof}

\section{T71: the local algebra at the zero potential}

Let $\zeta=e^{2\pi i/q}$ and introduce unnormalized Fourier coordinates
\begin{equation}\label{eq:fourier}
 h_k(v)=\sum_{j=0}^{q-1}v_j\zeta^{-kj}.
\end{equation}
On $\A_q$, direct reversal of the sum gives
\begin{equation}\label{eq:antiFourier}
 h_0=0,\qquad h_{-k}=-\zeta^{-k}h_k.
\end{equation}
Consequently $h_1,\ldots,h_m$ form a linear coordinate system on $\A_q$:
they determine all Fourier coefficients, and Fourier inversion determines $v$.

\begin{theorem}[The quadratic local algebra]\label{thm:origin}
For every $q\geq3$, the homogeneous quadratic parts of the $m$ equations
$f_{q,\ell}$ are an invertible linear combination of
$h_1^2,\ldots,h_m^2$. If $\mathfrak m$ is the maximal ideal at the origin,
then
\begin{equation}\label{eq:localgr}
 \gr_{\mathfrak m}\mathcal O_{\X_q,0}
 \simeq\C[h_1,\ldots,h_m]/(h_1^2,\ldots,h_m^2).
\end{equation}
Its Hilbert series is $(1+t)^m$, and
\begin{equation}\label{eq:originlength}
 \mu_q(0)=2^m.
\end{equation}
\end{theorem}

We first compute the quadratic terms explicitly. Write
$U_s=U_s(\lambda/2)$, where $U_0=1$, $U_1=\lambda$, and
$U_{s+1}=\lambda U_s-U_{s-1}$. Define cyclic correlations
\[
 C_d(v)=\sum_{j=0}^{q-1}v_jv_{j+d},\qquad
 K_d(\lambda)=U_{d-1}U_{q-d-1}\quad(1\leq d\leq m).
\]

\begin{lemma}[Quadratic discriminant]\label{lem:quadratic}
If $F_q^{[2]}$ denotes the homogeneous degree-two part in $u$, then
\begin{equation}\label{eq:quadratic}
 F_q^{[2]}(v;\lambda)=
 \sum_{1\leq d<q/2}C_d(v)K_d(\lambda)
 +\begin{cases}\frac12 C_m(v)K_m(\lambda),&q=2m,\\0,&q=2m+1.
 \end{cases}
\end{equation}
The $K_1,\ldots,K_m$ are linearly independent.
\end{lemma}
\begin{proof}
In the transfer product, replace the factors at positions $i<j$ by their
linear variations $-v_iE_{11}$ and $-v_jE_{11}$. With $A=A_0$ and $d=j-i$,
the resulting trace coefficient is
\[
 v_iv_j\tr\bigl(A^{q-d-1}E_{11}A^{d-1}E_{11}\bigr)
   =v_iv_jU_{q-d-1}U_{d-1},
\]
because $(A^s)_{11}=U_s$. Summing over unordered pairs gives
\eqref{eq:quadratic}. Antipodal pairs in even period are counted twice by
$C_m$, which explains the factor $1/2$.

The elementary Chebyshev product identity yields
\begin{equation}\label{eq:Kdifference}
 K_1=U_{q-2},\qquad K_d-K_{d-1}=U_{q-2d}\quad(2\leq d\leq m).
\end{equation}
For completeness, substitution $\lambda=2\cos\theta$ and
$U_s(\cos\theta)=\sin((s+1)\theta)/\sin\theta$ reduces the difference
to the sine subtraction formula. Polynomial identity then extends it to
all $\lambda$. The right sides of \eqref{eq:Kdifference} have the distinct
degrees $q-2,q-4,\ldots,q-2m$, proving independence.
\end{proof}

\begin{proof}[Proof of the quadratic assertion in Theorem~\ref{thm:origin}]
Lemma~\ref{lem:quadratic} says that the quadratic coefficient forms span
exactly the same space as $C_1,\ldots,C_m$. We now identify that space.
The correlations satisfy
\[
 C_{q-d}=C_d,\qquad
 \sum_{d=0}^{q-1}C_d=\left(\sum_jv_j\right)^2=0.
\]
Thus $C_0$ is a linear combination of $C_1,\ldots,C_m$.
Taking the discrete Fourier transform in $d$ gives
\begin{equation}\label{eq:correlationDFT}
 \sum_{d=0}^{q-1}C_d\zeta^{-kd}=h_kh_{-k}
        =-\zeta^{-k}h_k^2.
\end{equation}
Every $h_k^2$, $1\leq k\leq m$, therefore belongs to the span of
$C_1,\ldots,C_m$. Conversely, Fourier inversion in
\eqref{eq:correlationDFT}, using \eqref{eq:antiFourier}, expresses each
$C_d$ in the span of $h_1^2,\ldots,h_m^2$. Since these are squares of
independent coordinates, they are linearly independent quadratic forms.
Both spans have dimension $m$, so the change of generators is invertible.
\end{proof}

To pass from quadratic terms to the full local algebra, we record the precise
filtered-algebra fact being used.
\begin{lemma}[Regular lowest forms]\label{lem:initial}
Let $R=\C[[x_1,\ldots,x_m]]$ and let $g_1,\ldots,g_m\in R$ all have order
$d$. Suppose their lowest homogeneous forms $Q_1,\ldots,Q_m$ form a regular
sequence. Then
\[
 \gr\bigl(R/(g_1,\ldots,g_m)\bigr)
       =\C[x_1,\ldots,x_m]/(Q_1,\ldots,Q_m).
\]
\end{lemma}
\begin{proof}
The inclusion $(Q_1,\ldots,Q_m)\subseteq\ini(g_1,\ldots,g_m)$ is immediate.
For the reverse inclusion, take a nonzero $g=\sum_i a_ig_i$ and put
$n=\min_i\operatorname{ord}(a_i)$. If $n+d<\operatorname{ord}(g)$, the
degree-$n$ parts $A_i$ satisfy $\sum_iA_iQ_i=0$. First syzygies of a regular
sequence are the Koszul syzygies, so there are homogeneous $B_{ij}=-B_{ji}$
of degree $n-d$ such that $A_i=\sum_jB_{ij}Q_j$. Replacing
\[
 a_i\quad\hbox{by}\quad a_i-\sum_jB_{ij}g_j
\]
does not change $\sum_i a_ig_i$, since the $B_{ij}$ are skew symmetric,
and increases the minimum order of the coefficients. Repeat until
$n+d=\operatorname{ord}(g)$. This requires only finitely many steps. The
lowest part of $g$ now belongs to $(Q_1,\ldots,Q_m)$, as required.
The fact about syzygies is the first exactness statement of the Koszul
complex of a regular sequence. For an elementary induction, reduce a relation
modulo the ideal of the first $m-1$ forms. The last form is a non-zero-divisor
in that quotient, so its coefficient belongs to that ideal. Subtract the
corresponding Koszul relations and induct on the remaining $m-1$ forms.
\end{proof}

\begin{proof}[Completion of the proof of Theorem~\ref{thm:origin}]
An invertible constant change of the equations makes their lowest forms
$h_1^2,\ldots,h_m^2$. These form a regular sequence: after quotienting by
the preceding squares, multiplication by the next square remains injective
on the free module over the remaining polynomial variables.
Lemma~\ref{lem:initial} applies with $d=2$. Completing a finite local
$\C$-algebra does not change it, its associated graded algebra, or its
length. This proves \eqref{eq:localgr}. The square-free monomials
$\prod_i h_i^{\epsilon_i}$, $\epsilon_i\in\{0,1\}$, form a basis of the
graded quotient, giving Hilbert series $(1+t)^m$ and length $2^m$.
\end{proof}

\begin{remark}
The associated graded algebra is identified in \eqref{eq:localgr}; an
isomorphism of the original, ungraded local algebra with that square algebra
is not asserted. Moreover, existence in the next section needs only the
weaker upper bound $\mu_q(0)\leq2^m$, which follows already from containment
of the quadratic forms in the initial ideal.
\end{remark}

\section{T72: all-period antisymmetric existence}

\begin{theorem}[Conjecture~1 of \cite{FFORS}]\label{thm:existence}
For every integer $q>3$, there is a nonzero $v\in\A_q$ such that
$L_v(z)$ and $L_0(z)$ have identical characteristic polynomials for every
$z\ne0$. More precisely, with $m=\lfloor q/2\rfloor$,
\begin{equation}\label{eq:nonzeroLength}
 \sum_{v\in\X_q(\C),\,v\ne0}\mu_q(v)=2^m(m!-1).
\end{equation}
\end{theorem}
\begin{proof}
A finite algebra over $\C$ decomposes as the product of its local algebras.
By Proposition~\ref{prop:global}, the sum of their lengths is $2^m m!$.
Theorem~\ref{thm:origin} gives the contribution $2^m$ at zero. Subtracting
proves \eqref{eq:nonzeroLength}. When $q>3$, we have $m\geq2$, so its
right side is positive. Some nonzero point must therefore exist.
By \eqref{eq:coeffs}--\eqref{eq:scheme}, all its discriminant coefficient
differences vanish. Equation~\eqref{eq:floquet} gives the asserted
Floquet isospectrality, with no restriction on the period $q$.
\end{proof}

This argument does not interpret the global length as a count of distinct
points: it first removes the full local contribution of zero. That distinction
is essential. For $q=3$, $m=1$ and the total length equals the origin length,
so there is no nonzero antisymmetric point. Thus the threshold in the theorem
is sharp for $q\geq3$.

\section{T73: exact rank under period repetition}

For $w\in\A_r$ and $q=kr$, write $w^{[k]}$ for its concatenation $k$ times.
Reversal of a repeated word is the repetition of its reversal, so
$w^{[k]}\in\A_q$. Define a linear map on variations by
\begin{equation}\label{eq:S}
 S:\A_q\longrightarrow\A_r,\qquad
 (Sh)_j=\sum_{\ell=0}^{k-1}h_{j+\ell r}\quad(0\leq j<r).
\end{equation}
Reversal of the indices shows $(Sh)_{r-1-j}=-(Sh)_j$.
The map is onto, since $S(\eta^{[k]})=k\eta$.
Let $\Phi_q:\A_q\to\C^m$ be the coefficient map
$(f_{q,1},\ldots,f_{q,m})$, and write
$\rho_q(v)=\rank(d\Phi_q)_v$.

\begin{theorem}[The repeated-point differential]\label{thm:rank}
For $r\geq3$, $k\geq1$, $q=kr$, and any $w\in\A_r$,
\begin{align}
 D_q(w^{[k]};\lambda)&=2T_k\bigl(D_r(w;\lambda)/2\bigr),\label{eq:inflate}\\
 (dD_q)_{w^{[k]}}(h)&=
 U_{k-1}\bigl(D_r(w;\lambda)/2\bigr)(dD_r)_w(Sh),\label{eq:derivative}\\
 \rho_q(w^{[k]})&=\rho_r(w).\label{eq:rankEquality}
\end{align}
If $w\in\X_r(\C)$, then $w^{[k]}\in\X_q(\C)$ and
\begin{equation}\label{eq:tangent}
 \dim_\C T_{w^{[k]}}\X_q=\lfloor q/2\rfloor-\rho_r(w).
\end{equation}
For $k\geq2$, every such repeated point is singular.
\end{theorem}
\begin{proof}
The monodromy of a repeated potential is $M_r(w)^k$. For a determinant-one
$2\times2$ matrix $M$, Cayley--Hamilton implies
$\tr M^k=2T_k(\tr M/2)$, proving \eqref{eq:inflate}.

Let $\tau$ be cyclic shift by $r$ on $\C^q$ and let
$P=k^{-1}\sum_{a=0}^{k-1}\tau^a$. The point $w^{[k]}$ is fixed by $\tau$,
and the discriminant is invariant under $\tau$. Differentiating this
invariance gives
\[
 (dD_q)_{w^{[k]}}(h)=(dD_q)_{w^{[k]}}(Ph).
\]
The reversal operator $R$ satisfies $R\tau R=\tau^{-1}$, so $P$ commutes
with $R$ and preserves its negative eigenspace $\A_q$. Explicitly,
$Ph=(Sh/k)^{[k]}$. Differentiating \eqref{eq:inflate} along a repeated
variation $\eta^{[k]}$, and using $T_k'=kU_{k-1}$, gives
\[
 (dD_q)_{w^{[k]}}(\eta^{[k]})
   =kU_{k-1}(D_r(w)/2)(dD_r)_w(\eta).
\]
Substitution of $\eta=Sh/k$ proves \eqref{eq:derivative}.

The factor $U_{k-1}(D_r(w;\lambda)/2)$ is a nonzero polynomial: $D_r$ is
monic of degree $r$, and $U_{k-1}$ has degree $k-1$. Multiplication by it
is injective in $\C[\lambda]$. Since $S$ is onto, the images of the
differentials have equal dimension. Reading coefficients in the appropriate
parity polynomial spaces preserves that dimension, proving
\eqref{eq:rankEquality}.

If $w$ is zero-isospectral, the monodromy trace formula, or
$T_k(T_r(x))=T_{kr}(x)$, proves that its repetition is zero-isospectral.
The Zariski tangent space of the scheme \eqref{eq:scheme} is the kernel of
$d\Phi_q$, yielding \eqref{eq:tangent}. When $k\geq2$,
\[
 \rho_r(w)\leq\lfloor r/2\rfloor<\lfloor kr/2\rfloor.
\]
The tangent space is nonzero. A zero-dimensional local ring over $\C$ is
regular exactly when its tangent space is zero, so the point is singular.
\end{proof}

The trace identity \eqref{eq:inflate} is not claimed as new. The precise
content needed for the later applications is that the differential rank
does not increase when the period is enlarged, although the antisymmetric
ambient dimension does increase.

\section{T74: an exact odd-period counterexample}

We first give explicit regular seeds of periods four and five. Direct
transfer multiplication gives, for arbitrary $a,b\in\C$,
\begin{align}
 D_4(a,b,-b,-a)-D_4(0)
  &=-(a^2+b^2)\lambda^2+a^2b^2+a^2-2ab+b^2,\label{eq:seed4}\\
 D_5(a,b,0,-b,-a)-D_5(0)
  &=-(a^2+b^2)\lambda^3
     +(a^2b^2+2a^2-2ab+b^2)\lambda.\label{eq:seed5}
\end{align}
Here and below the suppressed spectral argument is $\lambda$.
Choose
\begin{equation}\label{eq:seeds}
 Z=(1+i,1-i,-1+i,-1-i),\qquad
 W=(a,ia,0,-ia,-a),\quad a^2=1-2i.
\end{equation}
Both choices of the square root $a$ work. Equation~\eqref{eq:seed4} vanishes
at $Z$: its quadratic coefficient is zero and its constant coefficient is
$4-4=0$. For $W$, $b=ia$ gives $a^2+b^2=0$, and the remaining coefficient is
\[
 -a^4+(1-2i)a^2=0.
\]
Thus $Z\in\X_4$ and $W\in\X_5$. Both seeds, up to cyclic shift and complex
conjugation, already occur in Papanicolaou \cite[Example~3(ii)--(iii),
pp.~807--809]{Papanicolaou}; the period-four solution also appears in
\cite{FLMPRTTZ}. Neither seed is claimed as a new discovery.

Taking the two nontrivial coefficients in the order displayed, their
Jacobian determinants with respect to $(a,b)$ are
\begin{equation}\label{eq:seedJacobians}
 \det(d\Phi_4)_Z=-16i,\qquad
 \det(d\Phi_5)_W=-16+12i.
\end{equation}
For example the first row is $(-2a,-2b)$; the second row for period five
is $(2ab^2+4a-2b,\,2a^2b-2a+2b)$. Substitution gives the second value in
\eqref{eq:seedJacobians}. Both seeds therefore have rank two.
The least period of $Z$ is four, since neither period one nor period two
is possible. The least period of the nonconstant word $W$ is five, since
five is prime.

\begin{theorem}[Odd singular points; refutation of Conjecture~3]\label{thm:odd}
With $W$ as in \eqref{eq:seeds}, for every odd $k\geq3$ the point
$W^{[k]}\in\X_{5k}$ is nonzero and singular. Its differential rank is two
and its Zariski tangent dimension is $(5k-1)/2-2$.
In particular, $\X_{15}$ has a nonzero singular point with tangent dimension
five. Hence the odd-period regularity assertion of Conjecture~3 in
\cite{FFORS} is false.
\end{theorem}
\begin{proof}
The seed identities and Theorem~\ref{thm:rank} prove every assertion.
Explicitly, at period fifteen the word is
\[
 (a,ia,0,-ia,-a,\ a,ia,0,-ia,-a,\ a,ia,0,-ia,-a).
\]
Its seven antisymmetric coordinates are
$(a,ia,0,-ia,-a,a,ia)$, so it lies in the prescribed specialization exactly.
Its discriminant is
\[
 D_{15}=D_5^3-3D_5,
 \qquad D_5=\lambda^5-5\lambda^3+5\lambda,
\]
which is the free discriminant. Its rank is $\rho_5(W)=2<7$.
The original statement quantifies over all $15\Z$-periodic potentials in
this specialization, including those of smaller least period. Thus this
example satisfies its hypotheses without changing the problem.
\end{proof}

\section{T75: even uniqueness fails, and singular classes are unbounded}

Use equivalence under cyclic shifts, reversal, global negation, and complex
conjugation, including all their compositions. This is at least as broad as
the symmetries used for the specialized scheme in \cite{FFORS}. Each of
these operations preserves least period.

\begin{theorem}[Multiple and unbounded singular classes]\label{thm:many}
The scheme $\X_{20}$ contains the nonzero singular points
\[
 P=Z^{[5]},\qquad Q=W^{[4]},
\]
where $Z,W$ are defined in \eqref{eq:seeds}. Both have differential rank
two and tangent dimension eight. They belong to different symmetry classes.
Consequently, the even-period uniqueness assertion of Conjecture~3 in
\cite{FFORS} is false.

More generally, for every integer $s\geq1$ there is an even period $q$
for which $\X_q$ has at least $s+1$ inequivalent nonzero singular points.
For every $s\geq2$ there is an odd period with at least $s$ such classes.
\end{theorem}
\begin{proof}
Theorem~\ref{thm:rank} and \eqref{eq:seedJacobians} give the membership,
rank, and tangent dimensions of $P,Q$. Their least periods are respectively
four and five. Since every permitted symmetry preserves least period, they
cannot be equivalent. This already disproves uniqueness at the allowed
even period $20>6$.

For the general assertion, choose distinct primes $p_1,\ldots,p_s\geq5$.
Theorem~\ref{thm:existence} supplies a nonzero $w_j\in\X_{p_j}$ for each
$j$. Such a word has least period $p_j$: the only smaller possible least
period divides the prime $p_j$ and hence is one, but a constant
antisymmetric word must be zero.
Let $q=4p_1\cdots p_s$. Repeat $Z$ and each $w_j$ to period $q$.
Their repetition factors are all greater than one, so all these points
are singular by Theorem~\ref{thm:rank}. Their distinct least periods
$4,p_1,\ldots,p_s$ make their symmetry classes distinct.
For the odd assertion, take $q=p_1\cdots p_s$ with $s\geq2$ and repeat
each $w_j$ to period $q$. The same argument applies.
\end{proof}

Theorems~\ref{thm:odd} and~\ref{thm:many} refute the two nonzero-point
predictions in Conjecture~3. Theorem~\ref{thm:origin} proves its remaining
origin-multiplicity assertion. No formula for the full local lengths of the
nonzero singular points is asserted here; tangent dimension alone does not
determine those lengths.

\section{Scope, prior work, and verification}\label{sec:scope}

The mathematical conclusions in this manuscript are unconditional over
$\C$. They depend on the displayed algebraic arguments, not on numerical
root finding or on an unproved conjecture. Independent internal derivations
checked the origin calculation both by cyclic correlations and by the
band-edge argument in Appendix~\ref{app:bandedge}. Exact symbolic checks of
the low-period seeds and repeated-point Jacobians provide additional tests;
they do not replace the general proofs. This manuscript has not undergone
external peer review or proof-assistant verification.

\begin{center}
\begin{tabular}{@{}p{.21\textwidth}p{.70\textwidth}@{}}
\toprule
Result & Precise contribution \\
\midrule
T71 & Full associated graded algebra and origin multiplicity, all $q\geq3$.\\
T72 & All-period affirmative solution of the existing Conjecture~1.\\
T73 & Exact differential-rank identity for repetition.\\
T74 & Period-15 counterexample and an infinite odd-period family.\\
T75 & Period-20 failure of uniqueness and unbounded singular classes.\\
\bottomrule
\end{tabular}
\end{center}

The prior global length calculation is explicitly credited in
Proposition~\ref{prop:global}. Papanicolaou's 2021 examples
\cite{Papanicolaou} supply both small-period seeds, and the 2024 construction
\cite{FLMPRTTZ} gives a broader even-period family.
The unrestricted existence theorem in \cite{CFK}
does not impose the reversal-antisymmetric constraint treated here, and
period inflation in that paper is a prior ingredient. The more recent
Jacobi-pencil work \cite{FaustJacobi} concerns reducibility and diagonal
equality patterns; no identical local-algebra identification was found in
the inspected text.

A literature check through 24 October 2026 examined these primary papers,
their available arXiv versions, associated author publication lists, and
searches for the exact conjecture numbers and the terms antisymmetric
Floquet isospectrality, spectral-ideal multiplicity, and repeated-potential
singularities. No earlier identical resolution was located in that search.
This is a bounded search result, not a proof that no such work exists
anywhere. In particular, this paper does not certify a historically
unprecedented foundational principle.

Against the requested success conditions: the targets are pre-existing
numbered conjectures and are used unchanged; the cited source presents
them as open; complete arguments are supplied here; the local computation
is not merely a renaming of the cited global count. The stronger priority
and historical-new-principle requirements remain uncertified. Therefore
this draft must not be advertised as having established every requested
historical condition. It also makes no claim to settle a transcendence
problem, Conjecture~2 of \cite{FFORS}, or five separate named open problems.

\appendix
\section{An independent band-edge proof of the quadratic structure}\label{app:bandedge}

This appendix gives a second proof of the structural assertion in T71.
It is useful for checking both the Fourier phase and the absence of a
degenerate quadratic direction.
For $1\leq k\leq m$, put
\[
 \theta_k=\frac{\pi k}{q},\qquad z_k=(-1)^k,\qquad
 \lambda_k=2\cos\theta_k,\qquad c_k=h_k/q.
\]
The real symmetric matrix $L_0(z_k)$ has the eigenvalue $\lambda_k$ with
multiplicity exactly two. An orthonormal basis of that eigenspace is
\[
 \varphi_+(j)=q^{-1/2}e^{i\theta_k j},\qquad
 \varphi_-(j)=q^{-1/2}e^{-i\theta_k j}.
\]
They obey the Floquet condition because $e^{\pm i\theta_k q}=z_k$.
They are distinct and orthogonal since $1\leq k<q$.
No other free eigenvector has the same eigenvalue, because equality of
the corresponding cosines forces one of these two angles modulo $2\pi$.

The compression of $\diag(v)$ to this eigenspace is
\[
 \begin{pmatrix}c_0&c_k\\c_{-k}&c_0\end{pmatrix}
       =\begin{pmatrix}0&c_k\\c_{-k}&0\end{pmatrix}.
\]
Decompose $\lambda_k I-L_v(z_k)$ into this eigenspace and its orthogonal
complement. Its complement block at $v=0$ is invertible. The Schur
complement therefore shows that the homogeneous quadratic term of its
determinant equals
\[
 \gamma_k(-c_kc_{-k})=\gamma_k\zeta^{-k}c_k^2,
\]
where $\gamma_k$ is the nonzero determinant of that complement block.
The off-diagonal coupling correction has order two, and replacing an
entry in the two-by-two leading block by that correction affects its
determinant only in degree at least three. This justifies taking exactly
the displayed quadratic term.

By \eqref{eq:floquet}, the determinant being expanded equals
$F_q(v;\lambda_k)$, since the free determinant vanishes at $\lambda_k$.
Evaluation at the $m$ values $\lambda_k$ is an invertible linear map on
the polynomial space in \eqref{eq:coeffs}. For even $q$, this is a
Vandermonde map in the distinct numbers $\lambda_k^2$, including zero.
For odd $q$, it is such a Vandermonde map followed by multiplication by
the nonzero $\lambda_k$. Thus the quadratic coefficient forms are an
invertible combination of the $m$ independent squares $c_k^2$. Applying
Lemma~\ref{lem:initial} gives the same local algebra and multiplicity.

\section{Exact computational checks and reproduction}

The accompanying verifier uses symbolic polynomial arithmetic, not floating
point approximations. Its checks include the following.
\begin{enumerate}
\item The discriminant identities \eqref{eq:seed4} and \eqref{eq:seed5} and
      the nonzero seed Jacobian determinants in \eqref{eq:seedJacobians}.
\item Vanishing of the discriminant differences at $W^{[3]}$, $Z^{[5]}$,
      and $W^{[4]}$, computed directly at periods 15 and 20.
\item Rank two of the full specialized coefficient Jacobians at those
      repeated points, using exact arithmetic in the seed fields.
\item For small periods, agreement of the quadratic transfer expansion
      with \eqref{eq:quadratic}, and length $2^m$ of the quotient by the
      homogeneous quadratic equations.
\end{enumerate}
For the period-five seed, exact reduction modulo $a^2-(1-2i)$ suffices;
for period four one may substitute $1+i$ directly or reduce modulo $a^2-2i$
with $b=-ia$. In the latter case both choices of $a$ give the stated ranks.
Each Jacobian column can be computed without expanding all potential
variables at once. If $B_j=A_{v_j}(\lambda)$ and $E=E_{11}$, then
\[
 \frac{\partial D_q}{\partial v_j}
   =-\tr\bigl(B_{q-1}\cdots B_{j+1}E B_{j-1}\cdots B_0\bigr).
\]
The column for antisymmetric coordinate $u_{j+1}$ is the difference
between the $v_j$ and $v_{q-1-j}$ derivative polynomials. Extracting the
coefficients $\lambda^{q-2\ell}$ produces the exact $m\times m$ Jacobian.
This is an independent computational check of the rank identity, rather
than a rank test on only repeated variations.

To compile this standalone source, run \texttt{pdflatex} twice on the
\texttt{.tex} file. No external figure or bibliography file is required.
The Python verifier and its machine-readable result are supplied as an
optional verification bundle. Its finite checks certify only the cases
executed; the all-period claims are established in the proofs above.

\begin{thebibliography}{9}
\bibitem{FFORS}
M. Faust, L. Friedman, G. O'Malley, R. Ramos, and A. Sharma,
\emph{Algebraic Properties of the Ideal of Spectral Invariants for the
Discrete Laplacian}, arXiv:2602.08304v1 (2026).
\url{https://arxiv.org/abs/2602.08304}.

\bibitem{CFK}
J. Cobb, M. Faust, and A. Kretschmer,
\emph{Inverse Eigenvalue Problems, Floquet Isospectrality and the
Hilbert--Chow Morphism}, arXiv:2601.09604v1 (2026).
\url{https://arxiv.org/abs/2601.09604}.

\bibitem{FLMPRTTZ}
M. Faust, W. Liu, R. Matos, J. Plute, J. Robinson, Y. Tao, E. Tran, and C. Zhuang,
\emph{Floquet Isospectrality of the Zero Potential for Discrete Periodic
Schr\"odinger Operators}, Journal of Mathematical Physics \textbf{65}
(2024), 073501. DOI: \href{https://doi.org/10.1063/5.0201744}{10.1063/5.0201744}.
\url{https://arxiv.org/abs/2401.09731}.

\bibitem{FaustJacobi}
M. Faust, \emph{Reducibility and Diagonal Equality Strata for Finite
Jacobi Pencils}, arXiv:2609.28490v1 (2026).
\url{https://arxiv.org/abs/2609.28490}.

\bibitem{Papanicolaou}
V. G. Papanicolaou,
\emph{Periodic Jacobi operators with complex coefficients},
Journal of Spectral Theory \textbf{11} (2021), no.~2, 781--819.
DOI: \href{https://doi.org/10.4171/JST/357}{10.4171/JST/357}.
\url{https://ems.press/journals/jst/articles/1865622}.
\end{thebibliography}
\end{document}
